\documentclass[10pt,a4paper]{amsart}
\usepackage[margin=19mm]{geometry}
\usepackage{amsmath,amssymb,mathtools}
\usepackage[scr=rsfs,cal=euler]{mathalfa}
\usepackage{xfrac}
\usepackage[unicode,hidelinks,bookmarks=false]{hyperref}
\newcommand{\ie}{i.e.\ }
\newcommand{\cf}{cf.\ }
\newcommand{\resp}{respectively\ }
\newcommand{\Subsection}{\subsection}
\providecommand{\nobreakdash}{\nobreak\hbox{-}}
\setlength{\emergencystretch}{2em}
\multlinegap0pt
\usepackage{amssymb}
\usepackage{dsfont}
\newtheorem{lemma}{Lemma}
\newtheorem{theorem}{Theorem}
\title{An analogue of the Gauss-Kuzmin problem for continued $A_{2}$-fractions}
\author{M.\ V.\ Pratsiovytyi, O.\ P.\ Makarchuk and S.\ P.\ Ratushniak}
\date{Source: arXiv:2606.03229v1 (CC BY 4.0)}
\begin{document}
\maketitle

\noindent\emph{Source and licence.} An analogue of the Gauss-Kuzmin problem for continued $A_{2}$-fractions. Version arXiv:2606.03229v1 (CC BY 4.0), distributed by arXiv under the Creative Commons Attribution 4.0 International licence, http://creativecommons.org/licenses/by/4.0/, as recorded in the arXiv licence metadata for this exact version and shown on the abstract page https://arxiv.org/abs/2606.03229v1. Full licence notice: \texttt{licenses/arxiv-2606.03229-CC-BY-4.0.txt}.\par\smallskip

\noindent\textit{Editorial scope.} Counted unit: the article's theorem mr (source0.tex lines 593--611). The article's complete original proof of it is delivered with it (source0.tex lines 612--669, one \texttt{proof} environment of 58 lines). No mathematics is written, repaired or re-derived: the statement and the proof are the article's own lines, in the article's own notation, and no retained line is substituted or re-flowed.  Retained uncounted as the source-local context the proof needs: lines 352--357; lines 441--448; lines 462--477; lines 484--486.  Nothing was omitted from any retained span, so the package carries no omission record.  No retained line contains a citation, so the wrapper carries no bibliography and no bibliography entry.  No retained line contains a figure environment, an image-inclusion command or drawing code, so no image is delivered and the package declares no asset dependency.  Editorial formatting only, in addition: an AMS \texttt{amsart} preamble that restates the article's own declarations and macros used by the retained lines, namely the environments \texttt{lemma}, \texttt{theorem} on separate counters, together with the standard packages those lines need.  The retained environments print in this document as Lemma 1, Theorem 1 in order of appearance, whereas the article prints the same items with the numbering of its own full document.  The complete native source of the article as distributed in the arXiv e-print for this exact version is bundled unchanged as \texttt{sources/201961/source0.tex}.
\begin{lemma}\label{lema3}
  For every $x\in I$ and every natural number $n$, the following inequalities hold:
  $$
  \frac{\frac{60}{7}}{(x+1)(x+2)}\leq f_n^\prime(x)\leq \frac{\frac{21}{2}}{(x+1)(x+2)}.
  $$
\end{lemma}

\begin{lemma}\label{lema6}
   If for some non-negative $t,T$, and natural $k$ one has $$
   \frac{t}{(1+x)(2+x)}\leq f_{k}^{\prime}(x)\leq\frac{T}{(1+x)(2+x)}\;\;\forall x\in I,
   $$ then
  \begin{equation}\label{eq}
    \frac{t}{(1+x)(2+x)}\leq f_{k+1}^{\prime}(x)\leq\frac{T}{(1+x)(2+x)}\;\;\forall x\in I.
  \end{equation}
 \end{lemma}

\begin{lemma}\label{lema7}
  If for some positive constants $C_1$, $C_2$ and a natural number $m$ 
  $$\frac{C_1}{(x+1)(x+2)}\leq   f_m^{\prime} (x) \leq  \frac{C_2}{(x+1)(x+2)} \;\; \forall x \in I,    $$
 then for every natural $k$
  $$\frac{g(C_1)}{(x+1)(x+2)}\leq   f_{m+k}^{\prime}(x) \leq  \frac{w(C_2)}{(x+1)(x+2)} \;\; \forall x \in I,    $$
  where
$$g(t)=t+\frac{15}{4}\left(1-t\ln\frac{10}{9}\right)-\left(\frac{99}{4}+\frac{384}{225}t\right)\cdot\gamma_k,$$
$$
w(t)=t-\frac{15}{4}\left(t\ln\frac{10}{9}-1\right)+\left(\frac{99}{4}+\frac{384}{225}t\right)\cdot\gamma_{k},
$$
and
\begin{equation}\label{maz}
\gamma_{n}=\frac{1}{\frac{4}{17}((\frac{5+\sqrt{17}}{4})(\frac{1+\sqrt{17}}{4})^{n}-(\frac{5-\sqrt{17}}{4})(\frac{1-\sqrt{17}}{4})^{n})
 ((\frac{3+\sqrt{17}}{2})(\frac{1+\sqrt{17}}{4})^{n}-(\frac{3-\sqrt{17}}{2})(\frac{1-\sqrt{17}}{4})^{n})}.
\end{equation}
\end{lemma}

\begin{equation}\label{eqit}
C_1 \leq \frac{1}{\ln\frac{10}{9}} \leq C_2.
\end{equation}

\begin{theorem}
For every $n\in N$ and $x\in I$, the following condition holds
\begin{equation}\label{mr}
  -\ln\frac{10}{9}|r_{[\sqrt{n}]}|\leq f_n(x)-\ln^{-1}\left(\frac{10}{9}\right)\ln\frac{5(x+1)}{3(x+2)}
\leq\ln\frac{10}{9}|d_{[\sqrt{n}]}|,
\end{equation}
where
$$
r_{n}=-\ln^{-1}\left(\frac{10}{9}\right)+\left(\frac{60}{7}-\frac{\frac{15}{4}-\frac{99}{4}\gamma_n}
{\frac{15}{4}\ln\frac{10}{9}+\frac{384}{225}\gamma_n}\right)\cdot \left(1-\frac{15}{4}\ln\frac{10}{9}-\frac{384}{225}\gamma_n\right)^{n-1}+
\frac{\frac{15}{4}-\frac{99}{4}\gamma_n}{\frac{15}{4}\ln\frac{10}{9}+\frac{384}{225}\gamma_n},
$$
$$
d_{n}=-\ln^{-1}\left(\frac{10}{9}\right)+\left(\frac{21}{2}-\frac{\frac{99}{4}\gamma_n+\frac{15}{4}}
{\frac{15}{4}\ln\frac{10}{9}-\frac{384}{225}\gamma_n}\right)\cdot \left(1-\frac{15}{4}\ln\frac{10}{9}+\frac{384}{225}\gamma_n\right)^{n-1}+
 \frac{\frac{99}{4}\gamma_{n}+\frac{15}{4}}{\frac{15}{4}\ln\frac{10}{9}-\frac{384}{225}\gamma_n},
$$
the sequence $(\gamma_{n})$ is defined by equality \eqref{maz},$[n]$ denotes the integer part of $n$.
\end{theorem}

\begin{proof}
Denote $g^{(n)}(x)=\underbrace{g(g(\ldots g(x)))}_{n},\quad w^{(n)}(x)=\underbrace{w(w(\ldots w(x)))}_{n}.
$
For a given natural number $t$, in Lemma~\ref{lema7} we take the pairs $(m,k)$ to be $(t,t)$,$(2t,t)$,\ldots,$((t-1)t,t).$

Taking into account Lemma \ref{lema3} and condition \eqref{eqit}, we successively obtain
$$
f_{2t}^\prime(x)\geq\frac{g(\frac{60}{7})}{(x+1)(x+2)} \quad \forall x\in I,
$$
$$
f_{3t}^\prime(x)\geq\frac{g^{(2)}(\frac{60}{7})}{(x+1)(x+2)} \quad \forall x\in I,
$$
$$
\ldots \quad \ldots \quad \ldots
$$
$$
f_{t^2}^\prime(x)\geq\frac{g^{(t-1)}(\frac{60}{7})}{(x+1)(x+2)} \quad \forall x\in I.
$$

It is easy to see that for the recursive sequence $x_{n+1}=a\cdot x_n+b$, with $a\neq1$, the following equality holds 
$$
x_n=\left(x_0-\frac{b}{1-a}\right)a^n+\frac{b}{1-a}.
$$
Taking $x_{0}=\frac{60}{7}, \quad a=1-\frac{15}{4}\ln\frac{10}{9}-\frac{384}{225}\gamma_t,  \quad b=\frac{15}{4}-\frac{99}{4}\gamma_t$, we obtain
$$
g^{(t-1)}(\frac{60}{7})=\left(\frac{60}{7}-\frac{\frac{15}{4}-\frac{99}{4}\gamma_t}{\frac{15}{4}
\ln\frac{10}{9}+\frac{384}{225}\gamma_t}\right)\cdot \left(1-\frac{15}{4}\ln\frac{10}{9}-\frac{384}{225}\gamma_t\right)^{t-1}+
\frac{\frac{15}{4}-\frac{99}{4}\gamma_t}{\frac{15}{4}\ln\frac{10}{9}+\frac{384}{225}\gamma_t}.
$$
Similarly,
$$
f_{t^2}^\prime(x)\leq\frac{w^{(t-1)}(\frac{21}{2})}{(x+1)(x+2)} \quad \forall x\in I,
$$
where
$$
w^{(t-1)}(\frac{21}{2})=\left(\frac{21}{2}-\frac{\frac{99}{4}\gamma_t+
\frac{15}{4}}{\frac{15}{4}\ln\frac{10}{9}-\frac{384}{225}\gamma_t}\right)\cdot \left(1-\frac{15}{4}\ln\frac{10}{9}+\frac{384}{225}\gamma_t\right)^{t-1}+
 \frac{\frac{99}{4}\gamma_{t}+\frac{15}{4}}{\frac{15}{4}\ln\frac{10}{9}-\frac{384}{225}\gamma_t}.
$$
Thus, for every natural number $t$
$$
\frac{\ln^{-1}\left(\frac{10}{9}\right)}{(x+1)(x+2)}+\frac{r_t}{(x+1)(x+2)}\leq f_{t^2}^\prime(x)\leq\frac{\ln^{-1}\left(\frac{10}{9}\right)}{(x+1)(x+2)}+\frac{d_t}{(x+1)(x+2)}.
$$

Taking $t=[\sqrt{n} ]$ and using Lemma~\eqref{lema6}, together with the fact that $[\sqrt{n}]^2\leq(\sqrt{n})^2=n$ we obtain
$$
\frac{\ln^{-1}\left(\frac{10}{9}\right)}{(x+1)(x+2)}+\frac{r_{[\sqrt{n}]}}{(x+1)(x+2)}\leq f_{n}^\prime(x)\leq\frac{\ln^{-1}\left(\frac{10}{9}\right)}{(x+1)(x+2)}+\frac{d_{[\sqrt{n}]}}{(x+1)(x+2)}.
$$
Denoting
$$f(x)=\frac{\ln\left(x+1\right)-\ln\left(x+2\right)+\ln\frac{5}{3}}{\ln\frac{10}{9}}=
\ln^{-1}\left(\frac{10}{9}\right)\ln\frac{5(x+1)}{3(x+2)}$$
and integrating the last inequality over the interval $[0,5;\tau]$ ($\tau \in I$), we obtain
 $$
\ln\left(\frac{10}{9}\right)r_{[\sqrt{n}]}f(\tau)\leq f_n(\tau)-\frac{\ln\left(\tau+1\right)-\ln\left(\tau+2\right)+\ln\frac{5}{3}}{\ln\frac{10}{9}}\leq\ln\left(\frac{10}{9}\right)
d_{[\sqrt{n}]}f(\tau),
$$
whence, taking into account that $0\leq f(x)\leq 1$ for  $x\in I$, equality~\eqref{mr} follows, which completes the proof.
\end{proof}

\end{document}
